\documentclass[11pt]{article}

\usepackage[margin=1in]{geometry}
\usepackage{amsmath,amssymb}
\usepackage{booktabs}
\usepackage{enumitem}
\usepackage{hyperref}
\hypersetup{hidelinks}

\title{Scaling Book Exercises: Section 12, Quiz 5}
\author{Lucas Yan}
\date{}

\begin{document}
\maketitle

\noindent
\textbf{Quiz:} LLM Rooflines\\
\textbf{Source scan:} \texttt{../handwritten/section\_12\_handwritten.pdf},
pages 11--13\\
\textbf{Reference:} \url{https://jax-ml.github.io/scaling-book/gpus/}

\section*{Question 1: B200 rooflines}

A B200 provides \(2250\) bf16 TFLOPs/s versus \(990\) for H100, a
\(2.25\times\) increase.  Within an eight-GPU node, egress bandwidth doubles
from \(450\) to \(900\) GB/s.  The B200 critical ratio is therefore
\[
\boxed{
\frac{2250\times10^{12}}{900\times10^9}
=2500.}
\]
For tensor parallelism the approximate limit is \(Y<F/2500\), close to the
H100 limit because compute and within-node communication improved together.

Beyond one node, egress remains \(400\) GB/s while compute increases.
The data-parallel critical per-GPU token batch becomes
\[
\boxed{
\frac{2250\times10^{12}}{400\times10^9}
=5625\text{ tokens/GPU}.}
\]
Thus B200 improves compute faster than scale-out bandwidth and makes the
cross-node communication roofline harder to satisfy.

\section*{Question 2: Sharding LLaMA 3-70B}

\subsection*{(a) Storage lower bound}

Training in bf16 with Adam requires approximately two bytes per parameter
for weights and eight bytes for fp32 optimizer state:
\[
70\times10^9(2+8)=\boxed{700\text{ GB}.}
\]
With 80 GB per H100, the raw lower bound is
\[
\left\lceil\frac{700}{80}\right\rceil
=\boxed{9\text{ GPUs}},
\]
or two eight-GPU nodes after rounding to the node topology.  This excludes
gradients and activation checkpoints.

\subsection*{(b) Training time}

The standard transformer estimate is six FLOPs per parameter per training
token:
\[
\text{total FLOPs}
=6(70\times10^9)(15\times10^{12})
=6.3\times10^{24}.
\]
At 45\% MFU on 4096 H100s, effective throughput is
\[
4096(990\times10^{12})(0.45)
\approx1.825\times10^{18}\text{ FLOPs/s}.
\]
Therefore
\[
\boxed{
T\approx\frac{6.3\times10^{24}}{1.825\times10^{18}}
=3.45\times10^6\text{ s}\approx40\text{ days}.}
\]

\subsection*{(c) Parallelism and communication}

For \(F=28{,}672\), the within-node H100 model-parallel roofline suggests
\[
\frac{F}{1995}\approx14.4.
\]
Because an H100 node contains eight GPUs, the practical choice is
\(\boxed{8}\)-way tensor parallelism.

Pure data parallelism would then require \(4096/8=512\) replicas, but each
GPU would hold
\[
\frac{700\text{ GB}}8=87.5\text{ GB},
\]
which exceeds its 80 GB HBM.

With ZeRO-3 plus eight-way tensor parallelism, memory fits, but the local
batch is only
\[
\frac{4\times10^6}{4096}\approx976\text{ tokens/GPU}.
\]
This is below the data-parallel roofline, and ZeRO-3 also communicates
parameters, so it remains communication-bound.

Adding eight-way pipeline parallelism makes each model shard span eight
nodes.  This raises the effective scale-out bandwidth to
\[
8(400)=3200\text{ GB/s},
\]
giving a critical batch of
\[
\boxed{
\frac{990\times10^{12}}{3200\times10^9}
\approx309\text{ tokens/GPU}.}
\]
The 976-token local batch is above this bound, so efficient pipelining can
make the configuration compute-bound.

\section*{Question 3: Megatron-LM configurations}

With sequence length \(4096\), the per-GPU token batches are
\begin{align*}
16\text{B:}\quad&
\frac{192(4096)}{192}
=\boxed{4096},\\
70\text{B:}\quad&
\frac{384(4096)}{768}
=\boxed{2048},\\
314\text{B:}\quad&
\frac{1536(4096)}{3072}
=\boxed{2048}.
\end{align*}

The unadjusted spine-level threshold is approximately
\(2472\) tokens/GPU.  The 16B model is already above it.  The 70B and 314B
models appear below it until their model parallelism is included:

\begin{center}
\begin{tabular}{@{}lrrl@{}}
\toprule
Model & Model parallelism & Adjusted threshold & Result \\
\midrule
16B  & --- & \(2472\) & compute-bound \\
70B  & \(16\times\) & \(\approx2472/2=1236\) & compute-bound \\
314B & \(64\times\) & \(\approx2472/8=309\) & compute-bound \\
\bottomrule
\end{tabular}
\end{center}

The \(16\times\) and \(64\times\) model-parallel layouts span approximately
two and eight nodes, improving effective spine throughput by \(2\times\)
and \(8\times\), respectively.  Hence all three configurations are
theoretically compute-bound under this roofline model.

\end{document}
